\documentclass[12pt,letterpaper,boxed]{hmcpset}
\usepackage[margin=1in,headheight=14pt]{geometry}
\usepackage{amsfonts, amsmath, amssymb, enumerate, fancyhdr, gensymb, lastpage, mathtools, parskip, graphicx}
\usepackage{xcolor, tikz-cd}
\newcommand{\wg}[1]{\textcolor{violet}{#1}}
\newcommand{\OO}{\mathcal O}
\newcommand{\Q}{\mathbb Q}
\newcommand{\R}{\mathbb R}
\newcommand{\C}{\mathbb C}
\newcommand{\Z}{\mathbb Z}
\newcommand{\abs}[1]{\left|#1\right|}
\newcommand{\im}{\text{im }}
\newcommand{\inv}{^{-1}}
\newcommand{\normal}{\unlhd} %% one can also use \trianglelelefteq
\newcommand{\anglee}[1]{\langle #1 \rangle}
\usepackage[shortlabels]{enumitem}

% Numbering macros
\pagestyle{fancy}
\lhead{Will Gilroy}
\chead{Algs Homework \#}
\rhead{03 November 2021}
\lfoot{}
\cfoot{}
\rfoot{Page\ \thepage\ of\ \pageref{LastPage}}

\linespread{1.5}

\newcommand\blankpage{
    \thispagestyle{empty}
    \addtocounter{page}{-1}
    \newpage}
\renewcommand\footrulewidth{0.4pt}

\begin{document}

\problemlist{Algorithms HW } 

%------------------------- Problem 1 -----------------------

\begin{problem}
	\includegraphics[scale=0.8]{1.png}
	\hfill
\end{problem}

\begin{solution}
Cosider the following limit, we have 
\begin{align*}
	\lim_{z \to 0}\abs{z f(z)}	
	&= \lim_{z \to 0}\abs{z}\abs{f(z)} \\
    &\leq \lim_{z \to 0}\abs{z}\frac{1}{\sqrt{\abs{z}}} \\
	&= \lim_{z \to 0}\sqrt{\abs z}\\
	&= 0,
\end{align*}
and so $\lim_{z \to 0} \abs{zf(z)} = 0$. 
I particular we then have that $\lim_{z \to 0} zf(z) = 0$ since $0$ is
the only complex number with zero modulus.

Now, we can construct a holomorphic extension of $f$ on $\Delta$.
(Note that this is the proof of an alternative characterization of 
the Riemann Removable singularity theorem that we saw in class).

Let $h: \Delta \to \C$ be defined by 
\[
	h(z) = \begin{cases}
		z^2f(z) & z \neq 0  \\
		0 & z = 0.
	\end{cases}
\]
Notice that $h$ is holomorphic on $\Delta^\star$ as the product of two
holomorphic functions on $\Delta^\star$. 
We show that $h$ is also holomorphic at $z = 0$. We have 
\[
	h'(0) = \lim_{z\to 0} \frac{z^2f(z) - 0}{z} 
	= \lim_{z \to 0} zf(z)  = 0,
\]
where the last equality follows by the argument above. And so $h$ is
holomorphic at $0$ and in particular $h'(0) = 0$.

Since $h$ is holomorphic at $z= 0$ it has a power series expansion
about $z = 0$
\[
	h(z) = c_0 + c_1 z + c_2 z^2+ c_3 z^3 + \cdots 
\]
However, by definition $c_0 = h(0) = 0$ and $c_1 = h'(0) = 0$ and so
we actually have 
\[
	h(z) = c_2 z^2 + c_3z^3 + \cdots 
\]
Recall that for $z \neq 0$ we have
\[
	f(z) = \frac{h(z)}{z^2} 
	= c_2 + c_3 z + c_4 z^2 + \cdots
\]
In other words, the power series
\[
	c_2 + c_3 z + c_4 z^2 + \cdots
\]
agrees with $f(z)$ for all $z \neq a$ in $\Delta^\star$. However, this power series ia
also manifestly holomorphic on all of $\Delta$. And so we have found a
holomorphic function which extends $f$ to $\Delta \to \C$. 


\end{solution}

\newpage

%------------------------- Problem 2 -----------------------

\begin{problem}[2]
	\hfill
\end{problem}

\begin{solution}
\end{solution}

\newpage

%------------------------- Problem 3 -----------------------

\begin{problem}
	
	\includegraphics[scale=0.8]{3.png}
	\hfill
\end{problem}
\begin{solution}
	Recall that the residue of a meromorphic function $f: U \to \C$ at
	a point $p \in U$ is given by the $a_{-1}$ coefficient in it's
	Laurent expansion about $p$.
	For this question we determine the Laurent expansion of
	$\sin^3(z)$ about $z = 0$ by using power series that we already
	know.

	Recall the Taylor expansion of $\sin z$ about $z= 0$ is given by \[
		\sin z = z - \frac{z^3}{3!}	+ \frac{z^5}{5!} - \cdots
	\]
	And so in particular we have the following Laurent series
	expansion of $\sin^3(1/z)$ about $z = 0$
	\[
	\sin^3\left(\frac{1}{z}\right)	= \left[
		\frac{1}{z} - \frac{1}{3!}\frac{1}{z^3} + \frac{1}{5!} \frac{1}{z^5} - \cdots
	\right]^3
	\]
	Now, notice that the largest power of $z$ in this expansion is
	going to be $1/z^3$ with all other terms having a lower power of
	$z$. That is the coefficients of the Laurent series expansion is
	$a_n = 0$ for all $n \geq -2$. In particular $a_{-1} = 0$ and so 
	the residue of $\sin^3(1/z)$ about $z = 0$ is $0$. 

	Note that althought the residue of our function is $0$ at $z=0$,
	our function does have a singularity at $0$. In fact, the Laurent
	series expansion about $z=0$ has infinitely many terms with negative powers of
	$z$ and so the singularity at $z= 0$ is essential.


\end{solution}

\newpage

%------------------------- Problem 4 -----------------------

\begin{problem}[4]
	\hfill
\end{problem}

\begin{solution}
\end{solution}

\newpage

\end{document}
